\subsection {\href{https://leetcode.com/problems/logical-or-of-two-binary-grids-represented-as-quad-trees/description/}{Logical OR of Two Quad-Trees (LeetCode 558)}}

\subsection*{Problem Description}
\vspace{-1.5em}
\noindent\rule{\textwidth}{0.4pt}

This problem asks us to merge two $n \times n$ binary matrices represented by Quad-Trees using a logical OR operation.\\
A Quad-Tree is a tree data structure where each internal node has exactly four children representing the four quadrants: top-left, top-right, bottom-left, and bottom-right. Each node in the tree has two properties:
\begin{itemize}
    \item \texttt{val}: A boolean value (True for 1s, False for 0s)
    \item \texttt{isLeaf}: A boolean indicating whether the node is a leaf (True) or has four children (False)
\end{itemize}

The construction follows these rules:
\begin{itemize}
    \item \textbf{If a region contains all the same values} (all 0s or all 1s):
        \begin{itemize}
            \item Create a leaf node with isLeaf = True
            \item Set val to the grid's value (True for 1, False for 0)
            \item The four children are set to null
        \end{itemize}
    \item \textbf{If a region contains mixed values} (both 0s and 1s):
        \begin{itemize}
            \item Create an internal node with isLeaf = False
            \item Set val to any value (it doesn't matter since it's not a leaf)
            \item Divide the current region into four equal quadrants
            \item Recursively construct four children nodes for each quadrant
        \end{itemize}
\end{itemize}

\subsection*{Solution Approach}
\vspace{-1.5em}
\noindent\rule{\textwidth}{0.4pt}

This solution using recursion to construct the Quad-Tree. Here is how the solution works step by step.\\
\textbf{Step 1: Base cases}
\begin{lstlisting}[language=C++,linewidth=\linewidth]
	if (!quadTree1 && !quadTree2) return nullptr;
	if (!quadTree1) return quadTree2;
	if (!quadTree2) return quadTree1;
	if (quadTree1->isLeaf && quadTree1->val) return quadTree1;
	if (quadTree2->isLeaf && quadTree2->val) return quadTree2;
	if (quadTree1->isLeaf && quadTree2->isLeaf) return quadTree1;
\end{lstlisting}
\begin{itemize}
    \item \textbf{Null checks:} If both tree is null, return null pointer. If either tree is null, the result simply inherits the remaining structure of the other tree.
    \item \textbf{Dominant All \texttt{1's} Leaves:} If either \texttt{quadTree1} or \texttt{quadTree2} is a leaf node with \texttt{val=1}, it dominates the OR operation. 
    \item \textbf{False Leaves:} If both nodes are leaves but both value is false, then return either tree.
\end{itemize}
\textbf{Step 2: Divide and merging}
\begin{itemize}
    \item If the base cases are not met, the current region must be subdivided.
    \item We continue call \texttt{intersect} and calculate on four quadrants (topLeft, topRight, bottomLeft, bottomRight)
    \item After the recursion returns from the four children, the current node evaluates them.
    \item If all four children are leaf nodes and they all share the exact same boolean value, the subdivision is redundant. Then we must convert the current internal node into a leaf node and removes the redundant child nodes.
\end{itemize}

\subsection*{Complexity Analysis}
\vspace{-1.5em}
\noindent\rule{\textwidth}{0.4pt}

\begin{itemize}
    \item \textbf{Time Complexity:} $\mathcal{O}(n^2 \log n)$ where $n$ is the side length of the grid.
    \item \textbf{Space Complexity:} $\mathcal{O}(n^2)$
\end{itemize}

\subsection*{Implementation} 
\vspace{-1.5em}
\noindent\rule{\textwidth}{0.4pt}
\lstinputlisting[language=C++,numbers=none]{code/leetcode558.tex}